\documentclass[a4paper,11pt]{article}
\usepackage[T2A]{fontenc}
\usepackage[utf8]{inputenc}
\usepackage[russian]{babel}
\usepackage{amsmath,amssymb,mathtools}
\usepackage{geometry}
\geometry{left=22mm,right=22mm,top=20mm,bottom=21mm}
\usepackage{microtype}
\usepackage{array,tabularx,booktabs}
\usepackage{parskip}
\usepackage{needspace}
\usepackage{xcolor}
\usepackage{tikz}
\usepackage{fancyhdr}
\usepackage[hidelinks]{hyperref}

\definecolor{accent}{HTML}{087B82}
\definecolor{darkaccent}{HTML}{075A60}
\definecolor{bodytext}{HTML}{282E31}
\definecolor{muted}{HTML}{667175}
\definecolor{errorred}{HTML}{B02A31}
\color{bodytext}
\setlength{\parindent}{0pt}
\setlength{\parskip}{8pt plus 1pt}
\setlength{\headheight}{22pt}
\renewcommand{\arraystretch}{1.32}
\pagestyle{fancy}
\fancyhf{}
\fancyhead[L]{\footnotesize\color{muted}Теория вероятностей. Независимые испытания}
\fancyhead[R]{\footnotesize\color{muted}Екатерина Гнедкова}
\renewcommand{\headrulewidth}{0.4pt}
\renewcommand{\headrule}{\hbox to\headwidth{\color{accent}\leaders\hrule height \headrulewidth\hfill}}
\fancyfoot[C]{\footnotesize\thepage}
\renewcommand{\footrulewidth}{0pt}
\newcommand{\reviewtitle}[1]{\section*{\color{darkaccent}#1}\vspace{-4pt}{\color{accent}\rule{\linewidth}{0.45pt}}\vspace{6pt}}
\newcommand{\subhead}[1]{\par\vspace{5pt}{\color{darkaccent}\bfseries #1}\par\vspace{-2pt}}
\newcommand{\wrong}{\textcolor{errorred}{\textbf{Ошибка:}}}
\newcommand{\idea}{\textcolor{darkaccent}{\textbf{Ключевая идея:}}}
\newcommand{\result}{\textcolor{darkaccent}{\textbf{Итог:}}}

\begin{document}
\pagenumbering{roman}
\thispagestyle{empty}
\begin{titlepage}
\centering
\vspace*{25mm}
{\large\color{darkaccent}МАТЕМАТИКА\par}
\vspace{24mm}
{\LARGE Ревью домашней работы\par}
\vspace{15mm}
{\Large\color{accent}Независимые испытания\par}
{\Large\color{accent}и формула Бернулли\par}
\vspace{18mm}
{\normalsize Подготовка к ЕГЭ по профильной математике\par}
\vfill
{\large Екатерина Гнедкова\par}
\vspace{11mm}
{\normalsize 8 ноября 2026 г.\par}
\vspace{19mm}
{\small Лёвин Артём Александрович\par}
{\small эксклюзивно для Екатерины Гнедковой\par}
\vspace{18mm}
\begin{tikzpicture}
\draw[accent,line width=0.9pt] (0,0) .. controls (1.5,0.65) and (3.1,-0.55) .. (4.5,0)
 .. controls (6.0,0.55) and (7.3,-0.5) .. (8.8,0);
\end{tikzpicture}
\end{titlepage}
\pagenumbering{arabic}

\reviewtitle{Сводная таблица}
\noindent\small
\begin{tabularx}{\linewidth}{@{}>{\raggedright\arraybackslash}p{12mm}>{\raggedright\arraybackslash}p{27mm}>{\raggedright\arraybackslash}X>{\raggedright\arraybackslash}p{29mm}>{\raggedright\arraybackslash}p{30mm}@{}}
\toprule
\textbf{№} & \textbf{Тема} & \textbf{Суть ошибки} & \textbf{Ваш ответ} & \textbf{Правильный ответ}\\
\midrule
\hyperref[task:4]{\textcolor{darkaccent}{\textbf{4}}} & Метод дополнения & Не учтены четыре исхода с ровно одним орлом & $\displaystyle\frac{15}{16}$ & $\displaystyle\frac{11}{16}$\\
\addlinespace[4pt]
\hyperref[task:7]{\textcolor{darkaccent}{\textbf{7}}} & Формула Бернулли & Вместо двух успехов подставлен один & $6\cdot\frac16\left(\frac56\right)^5$ & $\displaystyle\frac{3125}{15552}$\\
\addlinespace[4pt]
\hyperref[task:10]{\textcolor{darkaccent}{\textbf{10}}} & Метод дополнения & Не учтены шесть способов одной поломки & $\approx 0{,}37$ & $0{,}114265$\\
\bottomrule
\end{tabularx}
\normalsize
\vspace{10pt}
\textcolor{muted}{Нажмите на номер задания, чтобы перейти к его разбору.}

\clearpage
\thispagestyle{empty}
\vspace*{\fill}
\begin{center}
{\huge\color{darkaccent}Разбор ошибок}
\end{center}
\vspace*{\fill}
\clearpage


\phantomsection\label{task:4}
\reviewtitle{Задание 4. Орёл выпал более одного раза}

\subhead{Текст задания}
Правильную монету бросают четыре раза. Найдите вероятность того, что орёл выпадет более одного раза.

\subhead{Ваше решение}
В работе использована идея вычитания вероятности противоположного события из единицы. Итоговая запись читается как
\[
P=1-\frac{1}{16}=\frac{15}{16}.
\]
Для исключаемого события учтён один исход из шестнадцати. Значение $\frac{15}{16}$ получено арифметически верно из записанной разности, но сама вероятность исключаемого события найдена неправильно.

\subhead{Ошибка}
\wrong\ Неполно определено противоположное событие: при вычислении учтена лишь одна последовательность бросков вместо пяти.

\subhead{Где именно возникает ошибка}
Последний правильный шаг --- использование события, противоположного событию «орёл выпал более одного раза»:
\[
P(A)=1-P(\overline A).
\]
Верно и то, что при четырёх бросках правильной монеты имеется $2^4=16$ равновозможных исходов.

Первый неверный шаг --- запись $P(\overline A)=\frac{1}{16}$. Противоположное событие $\overline A$ означает, что орёл выпал \emph{не более одного раза}, то есть \emph{ни разу или ровно один раз}. Эти случаи вместе дают пять исходов.

\subhead{Почему это неверно}
Слова «более одного» означают два, три или четыре появления орла. Если событие $A$ не произошло, число орлов может быть нулём или единицей. При нуле орлов существует одна последовательность: все четыре раза выпала решка. При одном орле его место можно выбрать четырьмя способами: в первом, втором, третьем или четвёртом броске. Следовательно, противоположному событию соответствуют $1+4=5$ равновозможных исходов. Одного исхода недостаточно.

\subhead{Ключевая идея}
\idea\ Прежде чем вычитать вероятность противоположного события, перечислите все значения, при которых исходное условие не выполняется. Для события «более одного» противоположным является «ноль или один», а не один отдельно выбранный исход.

\Needspace{145pt}
\subhead{Исправление}
Сохраняем метод дополнения, но исправляем количество исключаемых исходов:
\[
P(\overline A)=\frac{1+4}{16}=\frac{5}{16}.
\]
Тогда
\[
P(A)=1-\frac{5}{16}=\frac{11}{16}.
\]

\Needspace{165pt}
\subhead{Полное правильное решение}
Обозначим через $A$ событие «орёл выпал более одного раза». Прямой подсчёт потребовал бы учитывать ровно два, три и четыре появления орла. Удобнее подсчитать исходы, в которых орёл появляется не более одного раза.

Каждый бросок имеет два равновероятных исхода, поэтому четыре броска дают
\[
N=2\cdot2\cdot2\cdot2=16
\]
равновозможных последовательностей.

Если орёл не выпал ни разу, все четыре броска дали решку: это ровно один исход. Если орёл выпал ровно один раз, его можно расположить на одной из четырёх позиций, значит таких исходов четыре. Эти два случая не пересекаются, поэтому число исходов противоположного события равно
\[
1+4=5.
\]
Получаем:
\begin{align*}
P(A)&=1-P(\overline A)\\
&=1-\frac{5}{16}\\
&=\frac{16-5}{16}=\boxed{\frac{11}{16}}=0{,}6875.
\end{align*}

\subhead{Проверка}
Можно посчитать требуемые исходы напрямую. Ровно два орла можно получить $\binom42=6$ способами, ровно три --- $\binom43=4$ способами, четыре орла --- одним способом. Всего
\[
6+4+1=11\quad\Longrightarrow\quad P(A)=\frac{11}{16}.
\]
Прямой подсчёт совпал с методом дополнения.

\subhead{Итог}
\result\ Противоположное событие содержит пять равновозможных исходов из шестнадцати. Вероятность того, что орёл выпал более одного раза, равна
\[
\boxed{\displaystyle\frac{11}{16}=0{,}6875}.
\]

\subhead{Задача на закрепление}
Правильную монету бросают пять раз. Найдите вероятность того, что орёл выпадет более одного раза. Ответ запишите в виде обыкновенной дроби.

\clearpage

\phantomsection\label{task:7}
\reviewtitle{Задание 7. Ровно два выпадения шестёрки}
\subhead{Текст задания}
Правильный игральный кубик бросают шесть раз. Найдите вероятность того, что шестёрка выпадет ровно дважды. Ответ запишите в виде выражения, затем вычислите дробь.

\subhead{Ваше решение}
В работе правильно записана общая формула Бернулли, однако при подстановке чисел получено выражение
\[
\frac{6!}{1!\,5!}\cdot\frac16\left(\frac56\right)^5
=6\cdot\frac16\left(\frac56\right)^5.
\]
Оно соответствует событию «шестёрка выпала ровно один раз». Числовая дробь для исходного события не приведена.

\subhead{Ошибка}
\wrong\ В формуле Бернулли взято неверное число успешных испытаний.

\subhead{Где именно возникает ошибка}
Последний правильный шаг --- выбор формулы для независимых испытаний:
\[
P(X=k)=\frac{n!}{k!(n-k)!}\,p^k(1-p)^{n-k}.
\]
Первый неверный шаг --- подстановка $k=1$. По условию требуется \emph{ровно два} выпадения шестёрки, поэтому $k=2$. Число остальных исходов тогда равно $n-k=4$.

\subhead{Почему это неверно}
В формуле $n$ --- общее число бросков, $k$ --- число бросков, в которых произошло интересующее нас событие, $p$ --- вероятность этого события в одном броске. Здесь $n=6$, $p=\frac16$, а успехом считается выпадение шестёрки. Слова «ровно дважды» однозначно задают $k=2$. Значение $k=1$ описывает другое событие: одна шестёрка и пять результатов, отличных от шести. Коэффициент $\frac{6!}{1!5!}$ и показатели степеней $1$ и $5$ тоже относятся к этому другому событию.

\subhead{Ключевая идея}
\idea\ Перед вычислениями выпишите три параметра: общее число испытаний $n$, требуемое число успехов $k$ и вероятность успеха в одном испытании $p$. Затем проверьте, что показатели степеней $k$ и $n-k$ в сумме дают $n$.

\subhead{Исправление}
Сохраняем правильно выбранную формулу и вероятность выпадения шестёрки. Меняем количество успехов и неуспехов:
\[
n=6,\qquad k=2,\qquad p=\frac16,\qquad 1-p=\frac56.
\]
Тогда
\[
P(X=2)=\frac{6!}{2!\,4!}\left(\frac16\right)^2\left(\frac56\right)^4.
\]

\Needspace{160pt}
\subhead{Полное правильное решение}
За шесть бросков необходимо получить две шестёрки и четыре результата, отличных от шести. Для одной фиксированной последовательности, например «шестёрка, не шестёрка, шестёрка, не шестёрка, не шестёрка, не шестёрка», вероятность равна
\[
\left(\frac16\right)^2\left(\frac56\right)^4.
\]
Нужно выбрать, какие два из шести бросков дадут шестёрку. Число таких вариантов:
\[
\binom62=\frac{6!}{2!\,4!}
=\frac{6\cdot5}{2\cdot1}=15.
\]
Поэтому искомая вероятность равна
\begin{align*}
P(X=2)&=15\left(\frac16\right)^2\left(\frac56\right)^4\\
&=15\cdot\frac1{36}\cdot\frac{625}{1296}
=\frac{9375}{46656}
=\boxed{\frac{3125}{15552}}.
\end{align*}

\subhead{Проверка}
Число благоприятных расположений равно $15$, поскольку порядок двух шестёрок в последовательности можно выбрать пятнадцатью способами. Дробь положительна и меньше единицы. Приближённо она равна $0{,}20094$, что согласуется с вероятностью события.

\subhead{Итог}
\result\ В формуле Бернулли здесь необходимо подставлять $k=2$. Ответ:
\[
\boxed{\displaystyle\frac{3125}{15552}}.
\]

\subhead{Задача на закрепление}
Правильную монету подбрасывают пять раз. Найдите вероятность того, что орёл выпадет ровно три раза. Ответ запишите в виде выражения и затем вычислите дробь.

\clearpage
\phantomsection\label{task:10}
\reviewtitle{Задание 10. Хотя бы два устройства}
\subhead{Текст задания}
Каждое из шести независимых устройств выходит из строя с вероятностью $0{,}1$. Найдите вероятность того, что сломаются хотя бы два устройства.

\subhead{Ваше решение}
Верно выбран переход к противоположному событию: из единицы следует вычесть вероятности того, что неисправностей нет и что сломалось ровно одно устройство. В работе записано:
\[
P(\text{хотя бы два})=1-0{,}1-0{,}531441\approx0{,}37.
\]
При этом вероятность отсутствия неисправностей вычислена верно: $0{,}9^6=0{,}531441$.

\subhead{Ошибка}
\wrong\ Вероятность выхода из строя ровно одного устройства принята равной $0{,}1$; не учтены другие пять возможных мест единственной неисправности.

\subhead{Где именно возникает ошибка}
Последний правильный шаг --- разложение события «хотя бы два неисправных устройства» с помощью противоположных случаев:
\[
P(X\ge2)=1-P(X=0)-P(X=1).
\]
Также правильно найдено $P(X=0)=0{,}9^6=0{,}531441$.

Первый неверный шаг --- запись $P(X=1)=0{,}1$. Число $0{,}1$ означает вероятность поломки \emph{одного заранее выбранного} устройства. Оно само по себе не описывает событие «сломалось ровно одно из шести».

\subhead{Почему это неверно}
Чтобы сломалось ровно одно устройство, требуется одновременно два условия: выбранное устройство сломалось, а остальные пять работают. Вероятность такого конкретного случая равна $0{,}1\cdot0{,}9^5$. Однако единственным неисправным может быть любое из шести устройств. Эти шесть случаев несовместны, поэтому их вероятности нужно сложить, то есть умножить на $6$.

\subhead{Ключевая идея}
\idea\ «Хотя бы два» удобно вычислять через два противоположных случая: ноль поломок и ровно одну поломку. Для \emph{ровно одной} поломки обязательно учитывайте число способов выбрать неисправное устройство.

\subhead{Исправление}
Сохраняем выбранный способ решения и исправляем один член:
\[
P(X=1)=\binom61(0{,}1)^1(0{,}9)^5
=6\cdot0{,}1\cdot0{,}9^5.
\]

\Needspace{160pt}
\subhead{Полное правильное решение}
Обозначим через $X$ количество неисправных устройств. Вероятность, что конкретное устройство исправно, составляет
\[
1-0{,}1=0{,}9.
\]
Все шесть устройств исправны, если каждое из них работает. В силу независимости
\[
P(X=0)=(0{,}9)^6=0{,}531441.
\]
Ровно одна поломка может произойти у любого из шести устройств. Для каждого выбора вероятность равна $0{,}1\cdot(0{,}9)^5$, поэтому
\begin{align*}
P(X=1)&=6\cdot0{,}1\cdot(0{,}9)^5\\
&=6\cdot0{,}1\cdot0{,}59049
=0{,}354294.
\end{align*}
Из единицы вычитаем вероятности обоих несовместных случаев:
\begin{align*}
P(X\ge2)&=1-P(X=0)-P(X=1)\\
&=1-0{,}531441-0{,}354294\\
&=\boxed{0{,}114265}.
\end{align*}

\subhead{Проверка}
Случаи «ни одной поломки», «ровно одна поломка» и «не менее двух поломок» образуют все возможные исходы. Сумма их вероятностей:
\[
0{,}531441+0{,}354294+0{,}114265=1.
\]
Следовательно, разбиение выполнено корректно. Вероятность не менее двух поломок составляет приблизительно $11{,}43\%$.

\subhead{Итог}
\result\ При поиске вероятности ровно одного успеха среди шести независимых испытаний требуется множитель $6$. Ответ:
\[
\boxed{0{,}114265}.
\]

\subhead{Задача на закрепление}
Пять независимых приборов выходят из строя с вероятностью $0{,}2$ каждый. Найдите вероятность того, что из строя выйдут хотя бы два прибора.

\end{document}